% Shared exercise chunk: l4-binary-count
\begin{exercisechunk}
\item \textbf{Repeated errors and their accumulation.}
\label[exercise]{ex:binary-count}
Compare isolated errors with errors that persist. Fix
\(T\ge2\), \(P=\delta_{0^T}\), \(0<\varepsilon<1\), and \(X_0=0\).
Use \(c_{\mathrm{cnt}}(x)=\sum_{t=1}^T x_t\) and
\(c_{\mathrm{fin}}(x)=x_T\).
The Markov generators \(Q_{\mathrm{ind}},Q_{\mathrm{bad}}\) use
\[
 \begin{array}{c|cc}
 &Q(X_{t+1}=1\mid X_t=0)&Q(X_{t+1}=1\mid X_t=1)\\ \hline
 Q_{\mathrm{ind}}&\varepsilon&\varepsilon\\
 Q_{\mathrm{bad}}&\varepsilon&1
 \end{array}
 \qquad(0\le t<T).
\]
For \(P\), use the kernel that returns zero from either state.
\begin{enumerate}[(a),beginpenalty=10000]
\item Show that both generators have conditional KL
\(-\log(1-\varepsilon)\) at every data prefix, and hence
\begin{equation}
 \KL(P, Q_{\mathrm{ind}})
 =\KL(P, Q_{\mathrm{bad}})
 =T\log\frac1{1-\varepsilon}.
 \label{eq:binary-kl}
\end{equation}
Compare this constant positionwise loss with the loss profile of
$Q_-$ and $Q_+$. Compute the probability that the first one occurs at
position \(s\), and the probability of no ones.
\item Prove
\begin{align}
 \E_{Q_{\mathrm{ind}}}[c_{\mathrm{cnt}}]&=T\varepsilon,
 \label{eq:ind-cost}\\
 \E_{Q_{\mathrm{bad}}}[c_{\mathrm{cnt}}]
 &=\sum_{t=1}^T[1-(1-\varepsilon)^t]
 =T-\frac{(1-\varepsilon)[1-(1-\varepsilon)^T]}{\varepsilon}.
 \label{eq:bad-cost}
\end{align}
Obtain the second identity both from marginal error probabilities
and by conditioning on the first error position. Also prove
\begin{equation}
 \E_{Q_{\mathrm{ind}}}[c_{\mathrm{fin}}]=\varepsilon,\qquad
 \E_{Q_{\mathrm{bad}}}[c_{\mathrm{fin}}]=1-(1-\varepsilon)^T.
 \label{eq:terminal-cost}
\end{equation}
\item Prove the finite bounds
\[
 \frac{\varepsilon T(T+1)}2
 -\frac{\varepsilon^2T(T+1)(T-1)}6
 \le\E_{Q_{\mathrm{bad}}}[c_{\mathrm{cnt}}]
 \le\frac{\varepsilon T(T+1)}2.
\]
Deduce \(\E_{Q_{\mathrm{bad}}}[c_{\mathrm{cnt}}]
\sim\varepsilon T(T+1)/2\) when \(T\varepsilon\to0\).
Explain why the expected count can grow quadratically with \(T\)
in this regime.
\item Show that both sequence TVs equal \(1-(1-\varepsilon)^T\).
For either generator, verify the bounds
\begin{align}
 \E_Q[c_{\mathrm{cnt}}]&\le T[1-(1-\varepsilon)^T],
 &&\text{using exact TV},\label{eq:binary-tv-bound}\\
 \E_Q[c_{\mathrm{cnt}}]&\le
 T\min\left\{1,\sqrt{\frac T2\log\frac1{1-\varepsilon}}\right\},
 &&\text{using Pinsker}.\label{eq:binary-pinsker-bound}
\end{align}
Compare these with the exact counts at
\(T=1000,\varepsilon=10^{-3}\), and their leading terms when
\(T\varepsilon\to0\). Which information about error duration
does the TV-to-count bound discard?
\item Set \(\varepsilon=1-e^{-\kappa/T}\) for fixed \(\kappa>0\).
Compute the two terminal costs and the ratio
\(\E_{Q_{\mathrm{bad}}}[c_{\mathrm{fin}}]/
\E_{Q_{\mathrm{ind}}}[c_{\mathrm{fin}}]\). Explain why this ratio is of order \(T\) even though the terminal cost range is
fixed. Which model attains the terminal-cost upper endpoint?
\end{enumerate}

\end{exercisechunk}
